
\clearpage
\begin{examplebox}
\textbf{五步路线·入门例 A｜先从初中圆的垂径定理出发。}
圆 $x^2+y^2=25$ 中，给弦的中点 $M=(3,2)$。
先检验 $3^2+2^2=13<25$，真实非退化弦存在。
由 $OM\perp AB$，有 $k_{OM}=2/3$，故
\[
\boxed{k_{AB}=-\frac32},\qquad
\boxed{y-2=-\frac32(x-3)}.
\]
\textbf{一句话：}在圆里，先想垂径定理，不必先设直线联立。
\textbf{再问``相加''：}由勾股定理，半弦长度
\[
MP=\sqrt{R^2-OM^2}=\sqrt{25-(3^2+2^2)}=2\sqrt3.
\]
这一步\textbf{不是再求一次方向}，而是在方向已经确定后，求要沿弦的方向走多远。
\par\medskip\noindent
\textbf{五步路线·入门例 B｜同一个中点，比较三条\textbf{真实存在}的弦。}
取 $M=(\frac12,1)$；椭圆、双曲线均取 $a=2,b=3$，圆取 $R=5$。
\[
\begin{array}{c|c|c}
\text{曲线}&\text{存在性检查}&k_{AB}\\\hline
x^2+y^2=25 &\frac14+1<25&-\frac12\\[3pt]
\frac{x^2}{4}+\frac{y^2}{9}=1 &\frac1{16}+\frac19=\frac{25}{144}<1&-\frac98\\[4pt]
\frac{x^2}{4}-\frac{y^2}{9}=1 & S=\frac1{16}-\frac19=-\frac7{144}<0&+\frac98
\end{array}
\]
\textbf{比较：}圆的斜率来自普通垂直；椭圆多出尺度 $b^2/a^2$；双曲线 $y^2$ 项变号后斜率符号相反。
三个结果都对应\textbf{实际存在的非退化弦}，不是只代公式算出的假弦。
\end{examplebox}
\clearpage
\begin{examplebox}
\textbf{新增三联例｜同一双曲线中，方向确定后，为什么还必须相加？}
统一采用原方程 $\mathcal H:\frac{x^2}{4}-y^2=1$，即 $a^2=4,b^2=1$。
\begin{enumerate}[leftmargin=2em,topsep=2pt,itemsep=3pt]
\item \textbf{$S<0$，确有异支弦：}给 $M=(0,1)$，则 $S=-1$。
相减可知 $u=4t,\ v=0$（水平），相加则给
$-1-4(-1)t^2=1$，故 $t^2=1/2$。
取 $t=1/\sqrt2$，得到
$P=(2\sqrt2,1)$、$Q=(-2\sqrt2,1)$，均在原双曲线上。
\item \textbf{$0<S<1$，有方向却无弦：}给 $M=(1,0)$，$S=1/4$。
相减得到 $u=0$（只能竖直）；相加给
$\frac14-v^2=1$，于是 $v^2=-\frac34$。
没有实位移，不存在真实端点。等价地，$t^2=(1/4-1)/(4\cdot1\cdot1/4)=-3/4$。
\item \textbf{$S>1$，确有同支弦：}给 $M=(3,0)$，$S=9/4$。
相减得 $u=0$（只能竖直）；相加给
$\frac94-v^2=1$，即 $v=\pm\sqrt5/2$。
端点 $P=(3,\sqrt5/2)$、$Q=(3,-\sqrt5/2)$ 均在右支。
\end{enumerate}
\textbf{一眼迁移：}这三例的差别不是``有没有候选方向''，而是
\textbf{沿唯一方向需要的半弦位移平方，究竟为正数还是负数}。
\end{examplebox}
\clearpage
\begin{examplebox}
\textbf{例 1｜会求斜率，不等于真有弦（最小反例）。}
已知 $\mathcal H:x^2/4-y^2/9=1$，问 $M=(2,1)$ 能否为一条非退化弦的中点？
\[
S=\frac{2^2}{4}-\frac{1^2}{9}=\frac89\in(0,1),
\quad k=\frac{9\cdot 2}{4\cdot 1}=\frac92.
\]
虽然候选线为 $y-1=\frac92(x-2)$，但 $0<S<1$，\textbf{没有两个实交点，故不存在。}
\par\medskip\noindent
\textbf{例 2｜异支与同支：同一条双曲线上，选中点一眼判断。}
对于 $\mathcal H:x^2/4-y^2/9=1$：
\begin{itemize}[leftmargin=2em,topsep=2pt,itemsep=1pt]
\item $M_1=(0,1)$：$S=-1/9$，候选线 $y=1$，端点
$A=(-2\sqrt{10}/3,1),\ B=(2\sqrt{10}/3,1)$，\textbf{异支}。
\item $M_2=(3,1)$：$S=77/36>1$，候选线 $y-1=\frac{27}{4}(x-3)$，\textbf{同支}。
\end{itemize}
\par\medskip\noindent
\end{examplebox}
\clearpage
\begin{examplebox}
\textbf{例 3｜固定斜率 $k=2$，求所有中点的轨迹（含取值范围）。}
\begin{center}\begin{tikzpicture}[x=.8cm,y=.8cm,font=\small,>=stealth]
\node[font=\small\bfseries] at (0,7.2) {固定弦斜率 $k=2$：中点轨迹是两条开射线};
\draw[gray,->] (-5.4,0)--(5.5,0) node[right] {$x_0$};
\draw[gray,->] (0,-5.8)--(0,5.8) node[above] {$y_0$};
\draw[gray,dashed,thick] (-3.023715784,-3.401680257)--(3.023715784,3.401680257);
\draw[teal,very thick,->] (3.023715784,3.401680257)--(4.9,5.5125);
\draw[teal,very thick,->] (-3.023715784,-3.401680257)--(-4.9,-5.5125);
\draw[teal,fill=white,thick] (3.023715784,3.401680257) circle(3pt);
\draw[teal,fill=white,thick] (-3.023715784,-3.401680257) circle(3pt);
\node[align=left] at (1.6,4.35) {$x_0=8/\sqrt7$\\$y_0=9/\sqrt7$};
\node[align=left] at (-1.6,-4.35) {$x_0=-8/\sqrt7$\\$y_0=-9/\sqrt7$};
\node[align=center] at (2.8,-2) {虚线段排除\\原点也排除};
\node[align=center] at (-2.8,1.8) {$y_0=9x_0/8$\\$|x_0|>8/\sqrt7$};
\node at (0,-6.45) {端点空心：$S=1$，两交点重合；双曲线为 $x^2/4-y^2/9=1$};
\end{tikzpicture}
\end{center}
在 $\mathcal H:x^2/4-y^2/9=1$ 中，$y_0\ne0$ 时
\[
2=\frac{9x_0}{4y_0}\ \Longrightarrow\ y_0=\frac98x_0,\qquad
S=\frac{x_0^2}{4}-\frac1{9}\left(\frac98x_0\right)^2
=\frac{7x_0^2}{64}.
\]
因为 $k=2>b/a=3/2$，原点处不存在斜率为 $2$ 的中心弦；且 $S\ge0$，必须 $S>1$，故
\[
\boxed{y=\frac98x,\qquad |x|>\frac8{\sqrt7}.}
\]
边界点仅对应切点重合，两条\textbf{开射线}才是完整的中点轨迹。
\par\medskip\noindent
\end{examplebox}
\clearpage
\begin{examplebox}
\textbf{例 4｜无斜率不等于无弦：竖直弦与中心弦。}
仍取 $x^2/4-y^2/9=1$。
\begin{itemize}[leftmargin=2em,topsep=2pt,itemsep=1pt]
\item $M=(3,0)$：$S=9/4>1$，竖直线 $x=3$ 与曲线交于 $(3,\pm 3\sqrt5/2)$，恰以 $M$ 为中点。
\item $M=(1,0)$：$S=1/4$，不能有非退化弦。
\item $M=O$：$S=0$，仍可取 $y=x$，因为 $|1|<3/2$；交点互为关于原点的对称点。
\end{itemize}
\par\medskip\noindent
\textbf{例 5｜综合变式：经过已知点且中点被限制在某条直线上。}
在 $x^2/4-y^2/9=1$ 上，是否存在一条弦过 $Q=(1,2)$，且中点在 $y=x$ 上？
若 $M=(m,m)$ 且 $m\ne0$，公式给 $k=9/4$。
过 $Q$ 的候选线 $9x-4y-1=0$ 与 $y=x$ 相交于
$M=(1/5,1/5)$，此时 $S=1/180\in(0,1)$，\textbf{排除}。
若 $M=O$，候选中心线只能是 $OQ$，即 $y=2x$，因 $2>3/2$ 无实交点。
其余竖直线过 $Q$ 只能为 $x=1$，也不与曲线相交。综上\textbf{无解}。
\par\medskip\noindent
\textbf{高价值变式（对照椭圆）。}
对于椭圆 $x^2/4+y^2/9=1$，给定 $M=(1,1)$，
候选斜率 $k=-9/4$；因为 $1/4+1/9=13/36<1$，实际存在一条非退化中点弦。
双曲线的斜率\textbf{符号相反}且\textbf{存在区域不能照搬椭圆内部判据}。
\end{examplebox}
